\documentclass[11pt, letterpaper]{article}

\usepackage[margin=1in, headheight=22pt, headsep=12pt]{geometry}
\usepackage{titlesec}
\usepackage{enumitem}
\usepackage{xcolor}
\usepackage{hyperref}
\usepackage{fancyhdr}
\usepackage{microtype}
\usepackage{parskip}
\usepackage{booktabs}
\usepackage{sectsty}

% --- Color Definitions ---
\definecolor{accent}{HTML}{2C3E6B}
\definecolor{subaccent}{HTML}{5B7BA5}
\definecolor{rulecolor}{HTML}{CCCCCC}
\definecolor{darktext}{HTML}{1A1A1A}

% --- Hyperlink Setup ---
\hypersetup{
    colorlinks=true,
    linkcolor=accent,
    urlcolor=subaccent,
    pdftitle={Early Blues and New Orleans Jazz},
    pdfauthor={Lecture Notes},
}

% --- Section Formatting ---
\titleformat{\section}
    {\Large\bfseries\color{accent}}
    {\thesection}{1em}{}
    [\vspace{-0.5em}\textcolor{rulecolor}{\rule{\textwidth}{0.4pt}}]

\titleformat{\subsection}
    {\large\bfseries\color{subaccent}}
    {\thesubsection}{1em}{}

\titlespacing*{\section}{0pt}{1.5em}{0.75em}
\titlespacing*{\subsection}{0pt}{1em}{0.5em}

% --- Header/Footer ---
\pagestyle{fancy}
\fancyhf{}
\renewcommand{\headrulewidth}{0.4pt}
\renewcommand{\headrule}{\hbox to\headwidth{\color{rulecolor}\leaders\hrule height \headrulewidth\hfill}}
\fancyhead[L]{\small\color{gray}Early Blues \& New Orleans Jazz}
\fancyhead[R]{\small\color{gray}Lecture Notes}
\fancyfoot[C]{\small\color{gray}\thepage}

% --- List Formatting ---
\setlist[itemize]{leftmargin=1.5em, itemsep=3pt, parsep=0pt, topsep=4pt}
\setlist[itemize,2]{leftmargin=1.5em, itemsep=2pt}

% --- Body Text ---
\color{darktext}

\begin{document}

% --- Title Block ---
\thispagestyle{empty}
\begin{center}
    \vspace*{1cm}
    {\Huge\bfseries\color{accent} Early Blues \&\\[0.3em] New Orleans Jazz}\\[1.5em]
    \textcolor{rulecolor}{\rule{0.6\textwidth}{0.8pt}}\\[1em]
    {\large Lecture Notes}\\[0.5em]
    {\color{gray}\small Evolution of Early Jazz from Blues Vocal Traditions\\through New Orleans Jazz and the Rise of the Soloist}
    \vspace{1.5cm}
\end{center}

% --- Table of Contents ---
\tableofcontents
\newpage

% ============================================================
\section{Blues Foundations}
% ============================================================

\begin{itemize}
    \item Blues derives from three main chords (I, IV, V) with additional chords often superimposed for harmonic richness
    \item The \textbf{blues scale} is a six-note pentatonic scale plus the flat 5 (tritone), showing connection to West African music
    \item The tritone was historically called the ``Devil's Interval'' due to its dissonant sound; it is central to Western music's tension-release structure
    \item \textbf{Blue notes}: flat 3, flat 5, flat 7 --- the characteristic intervals of the blues
    \item Standard format is \textbf{12-bar blues}, though earlier forms included 8-bar and 16-bar variations
    \item Recording technology limited early recordings to approximately three minutes per disc side
\end{itemize}

% ============================================================
\section{Key Technical Concepts}
% ============================================================

\subsection{Timbre}
The quality of sound that distinguishes different musicians and instruments. In jazz, having a unique, identifiable sound is a hallmark of greatness.

\subsection{Vibrato}
Oscillation or shakiness in sound; a vocal technique that heavily influenced instrumental jazz playing.

\subsection{Call and Response}
A fundamental structural pattern in blues and jazz. Early jazz features bends, moans, and blue notes in vocal delivery.

% ============================================================
\section{Bessie Smith --- ``Empress of the Blues''}
% ============================================================

\begin{itemize}
    \item Born 1894 in Tennessee; prominent figure in the early vocal blues tradition
    \item \textbf{``Backwater Blues''} (1927) is an iconic early jazz recording featuring James P. Johnson on piano
    \item Known for a darker timbre and extensive use of vibrato
    \item Her vocal techniques became models for instrumental jazz playing
\end{itemize}

% ============================================================
\section{Jelly Roll Morton}
% ============================================================

\begin{itemize}
    \item Born Ferdinand Morton in 1890 in New Orleans; died 1941. Creole composer and pianist
    \item Nickname ``Jelly Roll'' came from playing piano at houses of prostitution
    \item \textbf{``Dead Man Blues''} (1926), recorded with his \textit{Red Hot Peppers}, depicts a New Orleans funeral march
    \item Music was later criticized for sounding dated, stuck in a 1920s style
    \item Features clarinet and trumpet solos --- clarinet was prominent before the saxophone dominated in the 1930s--40s
    \item Recording has ragtime influence with stiff rhythm characteristic of early jazz
\end{itemize}

% ============================================================
\section{W.C. Handy --- ``Father of the Blues''}
% ============================================================

\begin{itemize}
    \item Pianist and trumpet player; important early composer
    \item \textbf{``St. Louis Blues''} (1922) was one of the first commercial jazz successes
    \item Featured a ``Spanish tinge'' (tango rhythm) in the composition
    \item Recording shows clear switches between swung and straight feel, minor and major tonalities
    \item Very little improvisation --- highly composed
    \item Notable use of muted trumpet changing the timbre
\end{itemize}

% ============================================================
\section{New Orleans Jazz Tradition}
% ============================================================

\begin{itemize}
    \item New Orleans is the birthplace of jazz, with instrumentation derived from marching bands
    \item \textbf{Typical ensemble}: trumpet/cornet, clarinet, trombone, banjo/guitar, drums, and tuba (before the upright bass became prominent)
    \item Features \textbf{collective improvisation} --- simultaneous melodic lines rather than solo-focused performance
    \item The \textbf{``second line'' rhythm} comes from the New Orleans funeral tradition: a somber procession followed by a celebration of life
    \item The traditional New Orleans jazz sound has been preserved to the present day
\end{itemize}

\subsection{Modern Example: Branford Marsalis}
\begin{itemize}
    \item \textbf{``Sidney in the House''} (1992) demonstrates a contemporary New Orleans jazz sound
    \item Features Branford Marsalis on soprano saxophone (paying homage to Sidney Bechet) and Wynton Marsalis on trumpet
    \item \textbf{Sidney Bechet} was important for making the soprano saxophone prominent in jazz before John Coltrane revived it in the 1960s
    \item Recording features second line rhythm transitioning to a contemporary swing feel
    \item \textbf{Preservation Hall} in New Orleans still performs in the traditional style
\end{itemize}

% ============================================================
\section{Louis Armstrong --- The First Great Soloist}
% ============================================================

\begin{itemize}
    \item Born 1901 in New Orleans; died 1971. Brought improvisation to the forefront as a soloist art form
    \item Self-taught orphan who aspired to a great sound from his first time playing trumpet
    \item Started with Joe ``King'' Oliver's Creole Jazz Band in 1923; emerged as a soloist with Fletcher Henderson's band, 1924--1925
    \item \textbf{Invented scat singing in 1926} with ``Heebie Jeebies''
    \item Trumpet playing featured a strident, bright sound with straight tone and vibrato at the end of notes --- a highly influential technique
    \item \textbf{``West End Blues''} (1928), from the Hot Five/Hot Seven recordings, is considered the first great improvised jazz solo
    \item The iconic opening cadenza is one of the most memorable in jazz history
    \item Nicknames: ``Satchmo'' and ``Pops''
\end{itemize}

\subsection{Armstrong's Complex Legacy}
\begin{itemize}
    \item Was as much entertainer as musician, drawing from the older tradition of minstrelsy
    \item Minstrel shows were among the most popular entertainment in 1910s--20s America
    \item Received criticism for a subservient persona, blackface performances, and stage antics
    \item Counter-argument: opened doors for Black artists as the first prominent crossover figure
    \item Harlem Renaissance figures called him the ``old Negro'' vs.\ Duke Ellington as the ``new Negro'' --- entertainment-focused vs.\ dignified
    \item Wynton Marsalis considers Armstrong his hero; Ken Burns's \textit{Jazz} documentary (2000) featured Armstrong prominently
    \item The documentary was controversial for emphasizing early jazz over modern developments (6--7 episodes on 1900--1950; final episode covering 1965--2000)
    \item Popular for singing (``What a Wonderful World,'' ``Hello, Dolly!'') but treasured in jazz for virtuoso trumpet playing
\end{itemize}

\subsection{Armstrong \& Fitzgerald Recording}
\begin{itemize}
    \item \textbf{``They Can't Take That Away from Me''} (1955) with Ella Fitzgerald demonstrates high-quality later work
    \item Shows Armstrong's trumpet playing and singing voice with better recording fidelity
    \item Features call and response between vocals and trumpet
    \item Armstrong's trumpet playing mirrors his vocal phrasing
    \item Longer structure than earlier recordings, based on popular song form
    \item More modern rhythm section with double bass and drums
\end{itemize}

% ============================================================
\section{Looking Ahead}
% ============================================================

Next lecture topics: \textbf{Ella Fitzgerald}, \textbf{Duke Ellington}, \textbf{Count Basie}, and the \textbf{Swing Era} of the 1930s.

\vfill
\begin{center}
    \textcolor{rulecolor}{\rule{0.4\textwidth}{0.4pt}}\\[0.5em]
    {\small\color{gray} Source: \href{https://www.notion.so/Early-Blues-and-New-Orleans-Jazz-Lecture-2f8ed208c188804fa22fdf626b96b492?pvs=189\#2f8ed208c18880a09061fb4270e21f96}{Notion --- Early Blues and New Orleans Jazz Lecture}}
\end{center}

\end{document}
